%% slides/11_ingestion.tex
\begin{frame}{Travaux Réalisés}{Pipeline d'Ingestion et Idempotence}
  \begin{center}
    \begin{tikzpicture}[node distance=6mm, scale=0.82, every node/.style={scale=0.82}]
      \node[node-gold] (start) {Données extraites brutes\\(\texttt{RawListingDraft})};
      \node[base-node, draw=ciDark, shape=diamond, aspect=2.2, below=6mm of start] (dec1) {L'URL existe en base ?};
      % Branche Gauche (Re-crawl)
      \node[node-white, left=1cm of dec1, text width=3.6cm] (recrawl) {\textbf{Déjà archivée}\\Mise à jour \texttt{last\_seen\_at}};
      \node[base-node, draw=ciDark, shape=diamond, aspect=2.2, below=6mm of recrawl] (dec2) {Le prix a changé ?};
      \node[node-green, below=6mm of dec2, text width=3.6cm] (hist-price) {Ajout d'un point de prix dans\\\texttt{listing\_history}};
      % Branche Droite (Nouveau)
      \node[node-white, right=1cm of dec1, text width=3.6cm] (new) {\textbf{Inconnue}\\Résolution adresse \& géographie};
      \node[node-green, below=8mm of new, text width=3.6cm] (match) {Traitement de déduplication DCS};
      \node[node-blue, below=8mm of match, text width=3.6cm] (insert) {Insertion PostgreSQL \& liaisons};

      \draw[flow-arrow] (start) -- (dec1);
      \draw[flow-arrow] (dec1.west) -- node[label-style, above] {Oui} (recrawl.east);
      \draw[flow-arrow] (dec1.east) -- node[label-style, above] {Non} (new.west);
      \draw[flow-arrow] (recrawl) -- (dec2);
      \draw[flow-arrow] (dec2.south) -- node[label-style, right] {Oui} (hist-price.north);
      \draw[flow-arrow] (dec2.west) -- node[label-style, above] {Non} ++(-1cm, 0) |- (hist-price.west);
      \draw[flow-arrow] (new) -- (match);
      \draw[flow-arrow] (match) -- (insert);
    \end{tikzpicture}
  \end{center}
\note{
\textbf{\textcolor{speakerKelcy}{[EN] AZEFACK K.}}\\
``When a crawler sends a listing, we first check whether we've already seen it. This is \textbf{idempotence}, and it prevents us from creating duplicates in our own database.''\\[6pt]
\textbf{\textcolor{speakerJacky}{[FR] MANGOUO J. — détaille le logigramme}}\\
``On extrait la donnée, on regarde si l'URL existe. Si \textbf{oui} — re-crawl — on met à jour la date de visite et on compare le prix : s'il a changé, on ajoute un point dans \texttt{listing\_history}. Si \textbf{non} — nouvelle annonce — on résout l'adresse, puis traitement DCS, puis insertion dans PostgreSQL.'' ~|~ \textbf{$\approx$1 min 30}
}
\end{frame}